\documentclass{article}
\usepackage{graphicx} % Required for inserting images
\usepackage{amsmath}
\usepackage{amsthm}
\usepackage{amssymb}
\usepackage{mathrsfs}
\usepackage{enumitem}
\newtheorem{theorem}{Theorem}
\newtheorem{lemma}{Lemma}
\newtheorem{corollary}{Corollary}
\theoremstyle{remark}\newtheorem{remark}{Remark}

% Theorem with manually specified number, to match Goldrei's numbering
\theoremstyle{plain}
\newtheorem{bookthminner}{Theorem}
\newenvironment{bookthm}[1]
  {\renewcommand\thebookthminner{#1}\bookthminner}
  {\endbookthminner}

\newcommand{\T}[0]{\mathcal{T}}
\newcommand{\B}[0]{\mathcal{B}}
\newcommand{\C}[0]{\mathcal{C}}
\newcommand{\R}[0]{\mathbb{R}}
\newcommand{\Q}[0]{\mathbb{Q}}
\newcommand{\Z}[0]{\mathbb{Z}}
\newcommand{\N}[0]{\mathbb{N}}
\newcommand{\e}[0]{\epsilon}


\begin{document}

\begin{bookthm}{2.4}
    The following are logical equivalences.
    \begin{enumerate}[label=\textup{(\alph*)}]
    \item $(\phi \to \psi) \equiv (\neg\phi \lor \psi) \equiv \neg(\phi \land \neg\psi) \equiv (\neg\psi \to \neg\phi)$
    \item $(\phi \iff \psi) \equiv ((\phi \to \psi) \land (\psi \to \phi))$
    \end{enumerate}
\end{bookthm}

\begin{proof}[Proof of (a)]
    We prove the three adjacent equivalences in turn and then chain them.

    \medskip\noindent
    \textbf{Claim 1:} $(\phi \to \psi) \equiv (\neg\phi \lor \psi)$.

    \smallskip\noindent
    Let $v$ be any truth assignment.

    \smallskip\noindent
    ($\Rightarrow$)\enspace
    Suppose $v((\phi \to \psi)) = F$. Then $v(\phi) = T$ and $v(\psi) = F$,
    so that $v(\neg\phi) = F$ and hence $v((\neg\phi \lor \psi)) = F$.

    \smallskip\noindent
    ($\Leftarrow$)\enspace
    Suppose $v((\neg\phi \lor \psi)) = F$. Then $v(\neg\phi) = F$ and
    $v(\psi) = F$, so that $v(\phi) = T$ and hence $v((\phi \to \psi)) = F$.

    \smallskip
    Since $v$ takes only the values $T$ and $F$, the contrapositive of the
    forwards direction proves
    \[
        v((\neg\phi \lor \psi)) = T \;\Rightarrow\; v((\phi \to \psi)) = T,
    \]
    while the contrapositive of the backwards direction proves
    \[
        v((\phi \to \psi)) = T \;\Rightarrow\; v((\neg\phi \lor \psi)) = T.
    \]
    Therefore $v((\phi \to \psi)) = F \iff v((\neg\phi \lor \psi)) = F$ and
    $v((\phi \to \psi)) = T \iff v((\neg\phi \lor \psi)) = T$, so that
    $v((\phi \to \psi)) = v((\neg\phi \lor \psi))$ for every truth
    assignment $v$. Thus $(\phi \to \psi) \equiv (\neg\phi \lor \psi)$.

    % Since implication is only false under one combination of truth values for $\phi$ and $\psi$, through bivalence it implicit that when the implication, $(\phi \to \psi)$, is true that $(\neg\phi \lor \psi)$ will also be true.

    % The truth table for $\lor$ demands both formulas to be false and so only yields that truth value under one possible truth assignment. Through bivalence it is implicit that when $(\neg\phi \lor \psi)$ is false that $(\phi \to \psi)$ will also be false.

    \medskip\noindent
    \textbf{Claim 2:} $(\neg\phi \lor \psi) \equiv \neg(\phi \land \neg\psi)$.

    \smallskip\noindent
    Theorem 2.3(h) provides us with the following De Morgan's Law:
    $\neg(\phi \land \psi) \equiv (\neg\phi \lor \neg\psi)$. Treating the
    formulas in Theorem 2.3(h) as metavariables and substituting
    $\phi \mapsto \phi$, $\psi \mapsto \neg\psi$ uniformly on both sides,
    we arrive at
    \[
        \neg(\phi \land \neg\psi) \equiv (\neg\phi \lor \neg\neg\psi).
    \]
    Now Theorem 2.3(g) states the law of double negation:
    $\neg\neg\phi \equiv \phi$. Substituting $\phi \mapsto \psi$ gives
    $\neg\neg\psi \equiv \psi$, so for every truth assignment $v$ we have
    $v(\neg\neg\psi) = v(\psi)$. By the truth table for $\lor$ it follows
    that $v((\neg\phi \lor \neg\neg\psi)) = v((\neg\phi \lor \psi))$ for
    every $v$, so that
    \[
        (\neg\phi \lor \neg\neg\psi) \equiv (\neg\phi \lor \psi).
    \]
    Chaining the two displayed equivalences by Exercise 2.35(c) and
    reversing the result by Exercise 2.35(b), we conclude that
    $(\neg\phi \lor \psi) \equiv \neg(\phi \land \neg\psi)$.

    \medskip\noindent
    \textbf{Claim 3:} $\neg(\phi \land \neg\psi) \equiv (\neg\psi \to \neg\phi)$.

    \smallskip\noindent
    Let $v$ be any truth assignment.

    \smallskip\noindent
    ($\Rightarrow$)\enspace
    Suppose $v(\neg(\phi \land \neg\psi)) = F$. Then
    $v((\phi \land \neg\psi)) = T$, so $v(\phi) = T$ and $v(\neg\psi) = T$,
    so that $v(\neg\phi) = F$ and hence $v((\neg\psi \to \neg\phi)) = F$.

    \smallskip\noindent
    ($\Leftarrow$)\enspace
    Suppose $v((\neg\psi \to \neg\phi)) = F$. Then $v(\neg\psi) = T$ and
    $v(\neg\phi) = F$, so that $v(\phi) = T$, hence
    $v((\phi \land \neg\psi)) = T$ and $v(\neg(\phi \land \neg\psi)) = F$.

    \smallskip
    Since $v$ takes only the values $T$ and $F$, the contrapositive of the
    forwards direction proves
    \[
        v((\neg\psi \to \neg\phi)) = T \;\Rightarrow\; v(\neg(\phi \land \neg\psi)) = T,
    \]
    while the contrapositive of the backwards direction proves
    \[
        v(\neg(\phi \land \neg\psi)) = T \;\Rightarrow\; v((\neg\psi \to \neg\phi)) = T.
    \]
    Hence $v(\neg(\phi \land \neg\psi)) = v((\neg\psi \to \neg\phi))$ for
    every truth assignment $v$, and so
    $\neg(\phi \land \neg\psi) \equiv (\neg\psi \to \neg\phi)$.

    \medskip\noindent
    Finally, since logical equivalence is transitive by Exercise 2.35(c),
    Claims 1--3 chain together to give
    \[
        (\phi \to \psi) \equiv (\neg\phi \lor \psi)
        \equiv \neg(\phi \land \neg\psi) \equiv (\neg\psi \to \neg\phi). \qedhere
    \]
\end{proof}

\begin{proof}[Proof of (b)]
    Let $v$ be any truth assignment.

    \smallskip\noindent
    (1)\enspace Let $v((\phi \iff \psi)) = T$. Then the truth table
    for $\iff$ tells us that $v(\phi) = T = v(\psi)$ or
    $v(\phi) = F = v(\psi)$. Suppose the former is true. Then
    $v((\phi \to \psi)) = T$ and $v((\psi \to \phi)) = T$, and so
    $v(((\phi \to \psi) \land (\psi \to \phi))) = T$. If instead
    $v(\phi) = F = v(\psi)$, then the antecedents of both implications are
    false, hence both implications take value $T$, and again
    $v(((\phi \to \psi) \land (\psi \to \phi))) = T$.

    \smallskip\noindent
    (2)\enspace Let $v((\phi \iff \psi)) = F$. Then the truth table
    for $\iff$ tells us that at most one of $\phi, \psi$ is true
    and at most one is false, so the two remaining possibilities are
    $v(\phi) = T, v(\psi) = F$ and $v(\phi) = F, v(\psi) = T$. Swapping
    $\phi$ and $\psi$ interchanges these two possibilities and, by the
    symmetry of the $\iff$ and $\land$ tables, changes neither
    side's value; so WLOG let $v(\phi) = T$ and $v(\psi) = F$. Then
    $v((\phi \to \psi)) = F$ and so
    $v(((\phi \to \psi) \land (\psi \to \phi))) = F$.

    \smallskip
    Since $v$ takes only the values $T$ and $F$, cases (1) and (2) are
    exhaustive. Hence
    $v((\phi \iff \psi)) = v(((\phi \to \psi) \land (\psi \to \phi)))$
    for every truth assignment $v$, so that
    $(\phi \iff \psi) \equiv ((\phi \to \psi) \land (\psi \to \phi))$.
\end{proof}


\end{document}
